% \section{Introduction}
% \label{sec:introduction}

% \IEEEPARstart{O}{ral} diseases represent a significant global health burden, affecting an estimated 3.5 billion individuals worldwide~\cite{ali2025interpretable}, with conditions such as dental caries, gingivitis, calculus, ulcers, hypodontia, and tooth discoloration accounting for the majority of dental visits. Early and accurate diagnosis is crucial for effective treatment planning and disease prevention. However, manual diagnosis from intraoral photographs faces substantial challenges including high inter-examiner variability (agreement rates as low as 65--75\% for certain conditions~\cite{peres2019oral}), frequent co-occurrence of multiple pathologies requiring simultaneous recognition, and real-world imaging variations that degrade reliability. Moreover, critical specialist shortages in underserved regions (ratios as low as 1:50,000 versus 1:2,000 in developed areas~\cite{speight2017screening}) underscore the urgent need for automated diagnostic systems to expand access to quality oral healthcare.

% Recent advances in deep learning have demonstrated remarkable capabilities in medical image analysis, with convolutional neural networks (CNNs) and Vision Transformers (ViTs) achieving expert-level performance across radiology, pathology, and ophthalmology~\cite{george2025deep}. In oral disease recognition specifically, recent benchmarks have achieved high accuracies on single-disease classification under controlled conditions, with EfficientNetB3 reaching 95.4\% (Kappa 0.989) on six oral disease classes~\cite{ali2025interpretable} and ConvNeXt achieving 81.06\% validation accuracy with strong F1-scores for isolated pathologies~\cite{george2025deep}. However, these results are obtained under idealized conditions with balanced datasets and consistent imaging protocols. Real-world deployment faces three critical challenges that remain inadequately addressed. First, \textit{severe class imbalance}: oral disease datasets exhibit extreme distribution skew, with prevalent conditions like gingivitis representing 24--31\% of samples while rarer conditions such as hypodontia constitute only 9--11\%, leading to majority-class bias where standard training achieves high overall accuracy by predicting only frequent classes while missing clinically important rare conditions~\cite{george2025deep,wang2024mosaic,song2021classification}. Second, \textit{disease co-occurrence}: clinical photographs frequently contain multiple concurrent pathologies (35--45\% of cases~\cite{araujo2025twostep}), yet most existing methods frame the problem as single-label classification~\cite{ali2025interpretable,lee2021classification}, failing to detect all present conditions simultaneously. Third, \textit{deployment robustness}: models trained on curated laboratory datasets exhibit substantial performance degradation when encountering real-world variations in lighting conditions, camera equipment, viewing angles, and image quality~\cite{araujo2025twostep,canali2022challenges}, raising concerns about reliability in diverse clinical settings.

% To address these challenges, we present a systematic framework for robust oral disease classification that emphasizes practical deployment considerations. Our approach implements a two-step detection-classification pipeline that first localizes candidate disease regions using object detection models (YOLOv8, DETR, Cascade R-CNN), then applies specialized classifiers to each extracted region. This architectural decomposition provides several advantages: (1) it naturally accommodates multiple concurrent diseases through independent region analysis, (2) enables targeted optimization strategies for each stage with appropriate loss functions, (3) allows higher-resolution analysis of individual regions compared to full-image processing, and (4) provides spatial interpretability through explicit region localization, enhancing clinical trust and enabling targeted treatment planning. Within this framework, we systematically address class imbalance through comprehensive training strategies combining Focal Loss for hard example emphasis, inverse-frequency class weighting, label smoothing for calibration, and strategic oversampling of minority classes. To ensure deployment reliability beyond standard accuracy metrics, we establish rigorous evaluation protocols including model calibration assessment through Expected Calibration Error (ECE), confidence distribution analysis, and systematic robustness testing under photometric perturbations, noise injection, blur effects, and compression artifacts that simulate real-world imaging variations.

% Our main contributions are threefold: \textbf{First}, we implement and systematically evaluate a two-step detection-classification framework on the challenging Oral Diseases dataset (12,817 images, 6 classes with natural imbalance reflecting clinical distributions), conducting comprehensive architectural exploration across 20+ state-of-the-art models including EfficientNet, RegNet, ResNet, DenseNet, and Vision Transformer families. Through systematic benchmarking, we identify optimal architectures achieving XXXX\% accuracy with F1-score of XXXX while providing clear deployment guidance for different resource constraints through efficiency analysis of inference speed and model size trade-offs. \textbf{Second}, we propose comprehensive imbalance-aware training strategies that synergistically combine Focal Loss, inverse-frequency class weighting, label smoothing, and SMOTE-inspired oversampling, demonstrating through ablation studies that this integrated approach achieves substantial improvements on minority classes compared to baseline cross-entropy training while maintaining overall performance. \textbf{Third}, we establish a reliability assessment framework that evaluates deployment readiness through multiple complementary metrics: calibration quality via Expected Calibration Error (ECE = XXXX), inter-rater agreement through macro-averaged Cohen's Kappa ($\kappa$ = XXXX), and robustness retention rates under systematic perturbations (brightness variations, Gaussian noise, blur, JPEG compression), demonstrating stable performance degradation curves suitable for clinical deployment. Through rigorous experimental validation, our framework provides practical guidance for deploying deep learning systems in real-world oral health screening applications where class imbalance, imaging variability, and reliability requirements pose significant challenges beyond laboratory benchmarks.


\section{Introduction}
\label{sec:introduction}

\IEEEPARstart{O}{ral} diseases represent a significant global health burden, affecting an estimated 3.5 billion individuals worldwide~\cite{ali2025interpretable}. 
% Lúc sau mention các loại bệnh vào sau
%, with conditions such as dental caries, gingivitis, calculus, ulcers, hypodontia, and tooth discoloration accounting for the majority of dental visits. 
Early and accurate diagnosis is crucial for effective treatment planning and disease prevention. However, manual diagnosis from intraoral photographs faces substantial challenges including high inter-examiner variability with agreement rates as low as 65\% to 75\% for certain conditions~\cite{peres2019oral}, frequent co-occurrence of multiple pathologies requiring simultaneous recognition, and real-world imaging variations that degrade reliability. Moreover, the need for automated diagnostic systems to expand access to quality oral healthcare underscores the importance of developing robust AI-assisted tools~\cite{speight2017screening}.

Recent advances in deep learning have demonstrated remarkable capabilities in medical image analysis, with convolutional neural networks (CNNs) and Vision Transformers (ViTs) achieving expert-level performance across radiology, pathology, and ophthalmology~\cite{george2025deep}. In oral disease recognition specifically, recent benchmarks have achieved high accuracies on single-disease classification under controlled conditions, with EfficientNetB3 reaching 95.4\% with Cohen's Kappa coefficient of 0.989 on six oral disease classes~\cite{ali2025interpretable} and ConvNeXt achieving 81.06\% validation accuracy with strong F1-scores for isolated pathologies~\cite{george2025deep}. However, existing approaches face critical limitations that hinder real-world clinical deployment. First, severe class imbalance: oral disease datasets exhibit extreme distribution skew, with prevalent conditions like gingivitis representing 24--31\% of samples while rarer conditions such as hypodontia constitute only 9--11\%, leading to majority-class bias in standard training~\cite{george2025deep,wang2024mosaic}. 
% Nhân nói đoạn này bị bias
% Second, severe class imbalance: oral disease datasets exhibit extreme imbalance, with some conditions overrepresented by a five- to ten-fold factor compared to rarer pathologies~\cite{george2025deep,wang2024mosaic}, leading to biased models that fail to detect minority classes~\cite{song2021classification,lee2021classification}.
Third, laboratory-to-clinic gap: models trained on curated datasets with consistent imaging protocols exhibit substantial performance degradation when deployed in real clinical environments with variable lighting, diverse camera equipment, and image quality variations~\cite{araujo2025twostep,canali2022challenges}.

% Chờ chốt methodology
%To address these challenges, we propose a \textit{two-step detection-classification framework} specifically designed for robust multi-label oral disease recognition in real-world conditions. Unlike conventional end-to-end approaches, our pipeline decomposes the problem into disease region localization followed by fine-grained classification of each detected region. This decomposition naturally handles multiple concurrent diseases, allows region-specific feature extraction at higher resolutions, and enables independent optimization of detection and classification objectives. However, implementing this framework for real-world deployment presents unique challenges including handling severe class imbalance, managing multi-label complexity with careful threshold selection, and maintaining robustness under diverse imaging conditions encountered in clinical practice.

To address these challenges, we develop a comprehensive framework for robust multi-class oral disease classification that tackles the limitations identified above through systematic architectural design and training strategies.

Our main contributions are threefold. First, we propose a \textit{lightweight hybrid Transformer–CNN architecture} for multi-class oral disease classification from intraoral images. While building upon established techniques (window attention from Swin Transformer, cross-attention fusion), we introduce specific adaptations: (i) a learnable gated residual connection with scaling parameter $\alpha$ for adaptive CNN-Transformer feature balancing, (ii) multi-scale cross-attention fusion at two hierarchical levels, and (iii) a compact design (15.7M parameters) that is 82\% smaller than ViT-Base (86M), 45\% smaller than Swin-Tiny (28.3M), and comparable to EfficientNet-B3 (12M) while incorporating global attention mechanisms absent in pure CNN architectures, enabling deployment on standard clinical hardware with inference time <50ms. Second, we develop \textit{comprehensive imbalance-aware training strategies} that systematically integrate Focal Loss, class-aware weighting, label smoothing, and oversampling techniques, resulting in substantial F1-score improvements, particularly for minority disease classes. Third, we introduce a \textit{reliability assessment framework} that evaluates model calibration using Expected Calibration Error (ECE) and assesses robustness under systematic input perturbations, providing quantitative measures of diagnostic reliability suitable for clinical deployment.